% ---------------------------------------------------------------------
%  Pages liminaires (numérotées en chiffres romains)
% ---------------------------------------------------------------------

\pagetitre{REMERCIEMENTS}

Nous tenons tout d'abord à remercier notre encadrante, Madame Hasnââ CHAABI, pour
la qualité de son accompagnement tout au long de ce projet. Ses conseils, sa
disponibilité et ses remarques nous ont aidés à garder une démarche rigoureuse, en
particulier sur l'évaluation des modèles et sur la façon de présenter nos résultats.

Nous remercions également l'École Marocaine des Sciences de l'Ingénieur (EMSI) pour la
formation reçue et pour les moyens mis à notre disposition durant notre cursus
d'ingénieur, en 5\textsuperscript{ème} année de la filière IASDR, option Intelligence
Artificielle. Nos remerciements vont aussi à l'ensemble des enseignants qui ont
contribué à notre formation, notamment en traitement automatique du langage, en
apprentissage profond, en génération de texte et en développement logiciel : ce projet
s'appuie directement sur ce qu'ils nous ont transmis.

Nous remercions enfin nos familles et nos proches pour leur soutien et leurs
encouragements, ainsi que les amis qui ont accepté de tester l'application et de nous
faire part de leurs retours. Ce projet a été pour nous l'occasion de construire, de bout
en bout, une application réelle qui associe développement mobile et intelligence
artificielle, et de mesurer ce qu'il faut pour faire passer un modèle de l'entraînement
à un usage quotidien.

% ---------------------------------------------------------------------
\pagetitre{RÉSUMÉ}

Ce projet porte sur la conception et le développement d'\equipay, une application
mobile de partage des dépenses de groupe (voyages, sorties, colocation, crédits) qui
intègre plusieurs modèles d'intelligence artificielle. L'objectif est double : permettre
à un groupe de savoir à tout moment qui doit combien à qui, et supprimer la saisie
fastidieuse qui décourage les utilisateurs des applications existantes.

L'application, développée avec Flutter, s'appuie sur Supabase (PostgreSQL, authentification,
stockage et temps réel) pour les données partagées, et sur un backend FastAPI exécuté
localement sur un Mac Apple Silicon pour les fonctions d'IA. Quatre usages de l'IA
générative y sont intégrés : un assistant conversationnel qui répond aux questions sur
l'application, un scanner de tickets de caisse fondé sur un modèle vision-langage, la
saisie d'une dépense à partir d'une simple phrase, et l'organisation automatique du
programme d'un voyage. Les trois premiers modèles sont des versions affinées par LoRA de
Qwen3-4B et Qwen3-VL-4B, entraînées sur GPU puis converties en 4 bits avec MLX pour
l'inférence ; le quatrième combine un calcul d'itinéraire (OSRM), un algorithme d'ordre de
visite et un modèle de langage chargé de rédiger le programme. Chaque entraînement est
suivi dans MLflow, et seule la meilleure version, désignée « champion », est servie par
l'application. Sur leurs jeux de test respectifs, l'assistant passe d'un score de 0,276 à
0,770, le scanner de 0,792 à 0,954 et la saisie par phrase de 0,656 à 0,976.

\noindent\textbf{Mots-clés :} Intelligence artificielle, Génération de texte, LLM,
Modèle vision-langage, LoRA, Quantification, MLX, MLOps, MLflow, Flutter, Supabase,
FastAPI, Partage de dépenses.

% ---------------------------------------------------------------------
\pagetitre{ABSTRACT}

This project covers the design and development of \equipay, a mobile application for
sharing group expenses (trips, outings, shared housing, loans) that embeds several
artificial intelligence models. The goal is twofold: let a group know at any time who
owes how much to whom, and remove the tedious data entry that discourages users of
existing apps.

The application is built with Flutter and relies on Supabase (PostgreSQL, authentication,
storage and real-time updates) for shared data, and on a FastAPI backend running locally
on an Apple Silicon Mac for AI features. Four generative AI use cases are integrated: a
conversational assistant that answers questions about the app, a receipt scanner based on
a vision-language model, expense entry from a single free-text sentence, and automatic
planning of a trip itinerary. The first three models are LoRA fine-tuned versions of
Qwen3-4B and Qwen3-VL-4B, trained on a GPU and then quantized to 4 bits with MLX for
inference; the fourth combines route computation (OSRM), a visit-ordering algorithm and a
language model that writes the final plan. Every training run is tracked in MLflow, and
only the best version, tagged as ``champion'', is served to the app. On their respective
test sets, the assistant improves from a score of 0.276 to 0.770, the receipt scanner from
0.792 to 0.954, and sentence-based entry from 0.656 to 0.976.

\noindent\textbf{Keywords:} Artificial Intelligence, Text Generation, LLM,
Vision-Language Model, LoRA, Quantization, MLX, MLOps, MLflow, Flutter, Supabase,
FastAPI, Expense Sharing.

% ---------------------------------------------------------------------
%  Table des matières, liste des figures, liste des tableaux
% ---------------------------------------------------------------------
\clearpage
\phantomsection
\addcontentsline{toc}{liminaire}{TABLE DES MATIÈRES}
{\setstretch{1.15}\tableofcontents}

\clearpage
\phantomsection
\addcontentsline{toc}{liminaire}{LISTE DES FIGURES}
{\setstretch{1.15}\listoffigures}

\clearpage
\phantomsection
\addcontentsline{toc}{liminaire}{LISTE DES TABLEAUX}
{\setstretch{1.15}\listoftables}

% ---------------------------------------------------------------------
\pagetitre{LISTE DES ABRÉVIATIONS}

\begin{table}[H]
\centering
\small
\begin{tabularx}{0.92\textwidth}{|L{2.6cm}|Y|}
\hline
\rowcolor{entete}\textbf{Abréviation} & \textbf{Signification} \\ \hline
IA      & Intelligence Artificielle \\ \hline
LLM     & \textit{Large Language Model}, grand modèle de langage \\ \hline
VLM     & \textit{Vision-Language Model}, modèle vision-langage \\ \hline
LoRA    & \textit{Low-Rank Adaptation}, adaptation par matrices de faible rang \\ \hline
PEFT    & \textit{Parameter-Efficient Fine-Tuning}, affinage économe en paramètres \\ \hline
MLX     & Bibliothèque d'apprentissage automatique d'Apple pour Apple Silicon \\ \hline
MLOps   & \textit{Machine Learning Operations}, industrialisation des modèles \\ \hline
GPU     & \textit{Graphics Processing Unit}, processeur graphique \\ \hline
API     & \textit{Application Programming Interface}, interface de programmation \\ \hline
REST    & \textit{Representational State Transfer} \\ \hline
JSON    & \textit{JavaScript Object Notation} \\ \hline
JWT     & \textit{JSON Web Token}, jeton d'authentification signé \\ \hline
RLS     & \textit{Row Level Security}, sécurité au niveau des lignes \\ \hline
SQL     & \textit{Structured Query Language} \\ \hline
OCR     & \textit{Optical Character Recognition}, reconnaissance de caractères \\ \hline
OSRM    & \textit{Open Source Routing Machine}, moteur de calcul d'itinéraires \\ \hline
UML     & \textit{Unified Modeling Language} \\ \hline
UI / UX & Interface utilisateur / expérience utilisateur \\ \hline
MAD, DH & Dirham marocain \\ \hline
\end{tabularx}
\caption{Liste des abréviations utilisées dans le rapport}
\end{table}
